\chapter{Hypotheses and Open Problems}
\label{ch:hypotheses}

Three hypotheses follow from the measurements. Each is stated with the evidence for it, a
prediction, and the observation that would refute it.

\section{H1. Friction tracks depth}
\label{sec:h1}

\textbf{Hypothesis.} The cost the protocol measures of a computation tracks the length of its
longest chain of dependent steps, and not the number of its steps.

\textbf{Evidence.} On the RTX 3070 an operation of no instruction costs 0 cycles a step, of one
instruction 4.05, and of two dependent instructions 8.05: cost grows by one instruction latency
per link of the chain. Two chains that share no link, interleaved a step at a time, cost the
dearer of the two alone in all 91 pairs measured. The interleaved cost is the larger of the two. In
scheduling this is the critical path \cite{src:Kelley-Walker}: the duration of a set of tasks is
the length of its longest dependent sequence, however many tasks run beside it.

\textbf{Reading.} The language takes friction as its measure of information
(Chapter~\ref{ch:language}). If H1 holds, what it measures is not the length of a shortest
description but something closer to Bennett's logical depth, the time a near-shortest program
takes to produce an object \cite{src:Bennett}. The two differ sharply. A random string has
maximal Kolmogorov complexity and almost no depth, and Bennett's own examples of shallow objects
are random sequences and perfect crystals. A computation at its Kolmogorov complexity can still
be cheap to run, and the protocol would read it as cheap. In measurable form, the principle
\emph{information is coherence} reads as \emph{cost is depth}.

\textbf{Prediction.} A dependent chain of \(d\) one-instruction operations costs about
\(4.05\,d\) cycles a step on this part, whatever independent work is interleaved with it, until
the interleaved work contends for a shared unit.

\textbf{Refutation.} Independent chains whose interleaved cost is the sum of their costs, or a
dependent chain cheaper than its depth times the instruction latency.

\section{H2. Bits repeat, costs vary}
\label{sec:h2}

\textbf{Hypothesis.} The states of the machine fall into two classes. States decided by which
sides held (lead, rite, void, wait, gray, blok) are reproducible exactly from run to run. States
decided by cost (busy against dual, tilt against drop, fuse against drop) vary with the host's load
and are reproducible only as rates. The gate and rank rule of Chapter~\ref{ch:protocol} holds
inside the state alphabet itself.

\textbf{Evidence.} Over 80 runs of one scenario every bit-decided state repeated in every run,
and every difference between runs was a both-held cycle reading busy in place of dual, at 5\% in
one batch and 9.4\% in the other (Section~\ref{sec:repeat}). Both sides of a pair came in past the
busy edge together 3.4 times as often as independence predicts. A slow moment of the host slows
both sides at once. busy reads load common to the pair, the common-mode noise of a measurement
that shares one host, and not a property of either address.

\textbf{Consequence.} A program should branch on bit-decided states and treat cost-decided states
as signals with a rate. busy should not be attributed to either address of a pair that shares a
host. Production monitoring draws the same line: saturation is a property of a shared resource,
and it is watched as a trend and not acted on per request \cite{src:SRE}.

\textbf{The spurious timeout.} By Proposition~\ref{prop:spurious} a timeout derived from \(k\)
exchangeable runs is beaten by a further run with probability below \(1/(k+1)\). In batch B the
timeouts were derived from 64 or 128 runs, a bound of 1.5\%. The measured rate over both batches
is 0.21\%, well inside it. Both late answers fell in runs whose timeouts were among the lowest,
derived in a quiet moment of the host. There the runs the timeout was derived from and the run
that beat it were not exchangeable: the host had changed between them. This is the tail latency
of large services \cite{src:Dean-Barroso} seen on one machine, and it is why a reference ask is
put beside each real one in place of a timeout carried forward.

\textbf{Prediction.} Two asks on separate hosts pass the busy edge together at about the product
of their separate rates.

\textbf{Refutation.} A bit-decided state that differs between runs whose recorded sides agree;
or, on separate hosts, joint passes far above the product.

\section{H3. A label carries its pair}
\label{sec:h3}

\textbf{Hypothesis.} A transition label loses nothing a branch can use where it is decided in its
situation and every passage through a superstate carries the state left and the state reached.

\textbf{Evidence.} Table~\ref{tab:labelbits}. The labels as printed keep 3.302 of 5.977 bits of
the transition. Deciding them in their situation keeps 3.848; writing sync as its passage, 4.856;
letting gray trips carry their state as well, 5.480. The 0.498 bits still lost sit in labels
that each name a rule covering many pairs, drop and halt the largest.

\textbf{Consequence.} The superstate form already used for sync extends to gray at no cost to the
atomicity of Proposition~\ref{prop:atomic}, since only the labels change and not where they fall.
The rule labels, drop, halt, join, tilt, vent and wake, can either carry their pair as sync does or
stay as names; that is a choice of the language's design and the measurement only prices it.

\textbf{Refutation.} A pair that cannot be recovered from a label carrying it, which would mean
two transitions share a pair and the machine is not a function of its pairs.

\section{Open problems}

\begin{enumerate}
\item \textbf{What the baseline is per.} A timeout is expressed against the baseline, and the
  baseline has no keying. Per host is coarse enough to carry signal and too coarse to remember a
  preference part by part. Per part and operation splits the samples until each is noise.
  Chaining clears the noise floor, and the chain length that makes a per-part baseline usable is
  not measured.
\item \textbf{Whether a handoff moves the channel.} Whether pass should make the right address
  the primary channel is undecided. In the field, failover usually does and failback is made a
  deliberate step.
\item \textbf{Recovery from collapse.} jamm, halt and fuse are named and have no defined
  recovery. The circuit breaker \cite{src:Nygard} gives one: stop asking for a waiting period,
  then let one trial ask through, closing on success and opening again on failure. In this
  language the trial succeeding is sprk.
\item \textbf{A loop construct.} The language has no syntax for a query loop. A loop has to
  express: a pair put each cycle; the timeout left implicit and derived, the unbound pass first;
  a branch on the label with its pair carried; a turn count that bounds the loop where a fault
  cannot reach it; and the record written as it runs. The field's nearest form is the control
  loop of a reconciling controller. The syntax is the language's to define.
\item \textbf{A difference inside the floor.} Whether a link whose difference between two
  branches sits inside the floor reads dual, both alike, or gray, neither known to be cheaper.
\item \textbf{Comparing branches.} The link-by-link comparison of two branches is stated and not
  yet measured.
\end{enumerate}

\section{Data and code availability}

The language and the programs that produce every figure in this research paper are in the orior
repository, \texttt{https://github.com/dstroy0/orior}. Each program writes the traces it reads.
